\BourbakiExerciseGroup

\begin{enumerate}
\def\labelenumi{\arabic{enumi})}
\tightlist
\item
  \begin{enumerate}
  \def\labelenumii{\alph{enumii})}
  \tightlist
  \item
    Let X be a topological space. Show that the following three statements are equivalent:
  \end{enumerate}
\end{enumerate}

\begin{enumerate}
\def\labelenumi{(\Alph{enumi})}
\setcounter{enumi}{16}
\tightlist
\item
  If \(x, y\) are any two distinct points of \(X\) , there is a neighbourhood of \(x\) which does not contain \(y\) .
\end{enumerate}

(Q') Every subset of X which consists of a single point is closed in X. (Q'\,') For every \(x \in X\) the intersection of the neighbourhoods of x consists of the point x alone.

A space X is said to be accessible if it satisfies these conditions.

\begin{enumerate}
\def\labelenumi{\alph{enumi})}
\setcounter{enumi}{1}
\tightlist
\item
  An ordered set X with the right topology (with respect to which X is a Kolmogoroff space {[}§ 1, Exercise 2){]} is accessible if and only if no two distinct elements of X are comparable, and then the right topology on X is the discrete topology.
\end{enumerate}

\begin{enumerate}
\def\labelenumi{\arabic{enumi})}
\setcounter{enumi}{1}
\tightlist
\item
  \begin{enumerate}
  \def\labelenumii{\alph{enumii})}
  \tightlist
  \item
    If a topology on a set X has a finite subbase, and if X is accessible with respect to this topology, then X is finite and the topology is discrete.
  \end{enumerate}
\end{enumerate}

\begin{enumerate}
\def\labelenumi{\alph{enumi})}
\setcounter{enumi}{1}
\tightlist
\item
  If, in an accessible space X, every intersection of open sets is open, then X is discrete (cf.~§ 7, Exercise 6).
\end{enumerate}

\begin{enumerate}
\def\labelenumi{\arabic{enumi})}
\setcounter{enumi}{2}
\tightlist
\item
  \begin{enumerate}
  \def\labelenumii{\alph{enumii})}
  \tightlist
  \item
    Let X be an accessible space, A a subset of X, and x a point of \(\overline{A}\) which does not belong to A. Show that every neighbourhood of x contains an infinite number of points of A.
  \end{enumerate}
\end{enumerate}

\begin{enumerate}
\def\labelenumi{\alph{enumi})}
\setcounter{enumi}{1}
\item
  Let X be an accessible space and let A be a subset of X. Show that the set of points \(x \in \overline{A}\) such that every neighbourhood of x contains a point of A other than x is closed in X.
\item
  Let X be an accessible space. Show that the intersection of all neighbourhoods of any subset A of X is equal to A.
\end{enumerate}

\begin{enumerate}
\def\labelenumi{\arabic{enumi})}
\setcounter{enumi}{3}
\tightlist
\item
  Show that every subspace of a Kolmogoroff space (resp. accessible space) is a Kolmogoroff space (resp. accessible space); and that a topology which is finer than a topology of a Kolmogoroff space (resp. accessible space) is the topology of a Kolmogoroff space (resp. accessible space).
\end{enumerate}

¶ 5) a) Let X be an infinite set. Show that all the topologies \(\mathcal{G}_{\mathfrak{c}}\) on X (§ 2, Exercise 8) such that \(1 < c \leqslant\) card (X) are topologies of an accessible space but are not Hausdorff.

\begin{enumerate}
\def\labelenumi{\alph{enumi})}
\setcounter{enumi}{1}
\tightlist
\item
  Show that the intersection of all the Hausdorff topologies on X is the non-Hausdorff topology \(\mathfrak{G}_{\mathrm{card}(\mathbf{N})}\) , which is also the coarsest topology of an accessible space on X. (Use Exercise 8 of §2, and remark that there exist Hausdorff topologies on X in which there are countably infinite subsets of X which are not closed).
\end{enumerate}

\begin{enumerate}
\def\labelenumi{\arabic{enumi})}
\setcounter{enumi}{5}
\tightlist
\item
  \begin{enumerate}
  \def\labelenumii{\alph{enumii})}
  \tightlist
  \item
    Let X be a Hausdorff space and D a dense subset of X. Show that \(\text{card}(X) \leqslant 2^{2^{\text{card}(D)}}\) (observe that every point of X is the limit of a filter base on D). Show that this upper bound for card (X) cannot be improved (§ 4, Exercise 5 b)) and deduce that on an infinite set A the set of ultrafilters and the set of filters are each equipotent to \(\mathfrak{P}(\mathfrak{P}(A))\) .
  \end{enumerate}
\end{enumerate}

\begin{enumerate}
\def\labelenumi{\alph{enumi})}
\setcounter{enumi}{1}
\tightlist
\item
  Let K be the discrete space consisting of the two numbers o and 1. Let A be an infinite set. Deduce from a) that if card (A) ≥ 2 \(^{2\text{card(N)}}\) , then the product space K \(^{\text{A}}\) satisfies the condition (D \(_{\text{IV}}\) ) of Exercise 7 of § 1, but satisfies neither (D \(_{\text{II}}\) ) nor (D \(_{\text{III}}\) ). (§ 4, Exercise 4).
\end{enumerate}

* 7) Let X be the interval \([-1, 1]\) in R. Consider the following equivalence relation S on X: if \(x \neq \pm 1\) , the equivalence class of \(x\) consists of \(x\) and \(-x\) ; the equivalence class of \(1\) (resp. \(-1\) ) consists of \(1\) (resp. \(-1\) ) alone. Show that S is open and that the quotient space X/S is accessible but not Hausdorff. Show also that every point of X/S has a neighbourhood in X/S which is a regular subspace. *

\begin{enumerate}
\def\labelenumi{\arabic{enumi})}
\setcounter{enumi}{7}
\tightlist
\item
  \begin{enumerate}
  \def\labelenumii{\alph{enumii})}
  \tightlist
  \item
    Let X be a Hausdorff (resp. regular) space which is the sum of a family \((\mathbf{X}_{t})\) of subspaces, and let X/R be the quotient space obtained by pasting together the \(X_{t}\) along closed subsets \(A_{tx}\) by means of homeomorphisms \(h_{x_{t}}\) (§ 3, no. 4); suppose also that for each index t there is only a finite number of indices x such that \(A_{tx} \neq \emptyset\) . Under these conditions, show that X/R is Hausdorff (resp. regular).
  \end{enumerate}
\end{enumerate}

* b) Show that the space X/S of Exercise 7 can be obtained by pasting together two regular subspaces along two open sets. *

\begin{enumerate}
\def\labelenumi{\arabic{enumi})}
\setcounter{enumi}{8}
\tightlist
\item
  Let \(\Gamma\) be a finite group of homeomorphisms of a Hausdorff (resp. regular) space X, and let R be the equivalence relation ``there exists \(\sigma \in \Gamma\) such that \(y = \sigma(x)\)'' between \(x\) and \(y\) in X. Show that the quotient space X/R is Hausdorff (resp. regular).
\end{enumerate}

* 10) Let X be the subspace of R which is the complement of the set of points \(\mathbf{I} / n\) , where \(n\) is a rational integer other than \(\mathbf{o}, \mathbf{i}\) or \(-\mathbf{i}\) . Let S be the equivalence relation on X whose equivalence classes are (i) the set \(\{\mathbf{o}\}\) ; (ii) the set of all non-zero integers; (iii) all sets of the form \(\{x, \mathbf{i} / x\}\) where \(x \in \mathbf{X}\) and \(\mathbf{o} < |x| < \mathbf{i}\) .

\begin{enumerate}
\def\labelenumi{\alph{enumi})}
\item
  Show that the relation S on X is neither open nor closed.
\item
  Show that the graph of S is closed in \(X \times X\) , but that the quotient space X/S is not Hausdorff. *
\end{enumerate}

¶ * 11) Let X be a topological space. Let Z denote the product space \(R^{X}\) , and for each \(x \in X\) let \(Z_{x}\) be the subset of Z consisting of all \(u \in R^{X}\) such that \(u(x) \in Q\) and \(u(y) \notin Q\) for all \(y \in X\) other than x. Let Y be the subspace of the product \(X \times Z\) which is the union of the sets \(\{x\} \times Z_x\) as \(x\) runs through X. Show that Y is Hausdorff and that X is homeomorphic to a quotient space of Y. *

\begin{enumerate}
\def\labelenumi{\arabic{enumi})}
\setcounter{enumi}{11}
\tightlist
\item
  Let X be a topological space which has at least two points.
\end{enumerate}

\begin{enumerate}
\def\labelenumi{\alph{enumi})}
\item
  Show that the topologies \(\mathfrak{C}_{\Omega}\) and \(\mathfrak{C}_{\Phi}\) on \(\mathfrak{P}_0(\mathbf{X})\) (§ 2, Exercise 7) are not topologies of an accessible space.
\item
  Let \(\mathcal{G}_{\Theta}\) denote the least upper bound of the topologies \(\mathcal{G}_{\Omega}\) and \(\mathcal{G}_{\Phi}\) . Show that the set \(\mathfrak{F}(\mathbf{X})\) of non-empty closed subsets of \(\mathbf{X}\) , with the topology induced by \(\mathcal{G}_{\Theta}\) , is a Kolmogoroff space, and that if \(\mathbf{X}\) is accessible, so is \(\mathfrak{F}(\mathbf{X})\) .
\item
  Supposing that X is accessible, show that the space \(\mathfrak{F}(\mathbf{X})\) (with the topology induced by \(C_{\Theta}\) ) is Hausdorff if and only if X is regular.
\end{enumerate}

¶ 13) Let X be a set and let Γ be a group of permutations of X. a) Let F be a set of subsets of X. Let S be the set of all topologies on X in which all the sets of F are closed and every u ∈ Γ is a homeomorphism. Show that the set S has a least element C(Γ, F). If X is infinite, show that there exist sets F of subsets of X such that C(Γ, F) is accessible but not discrete (§ 2, Exercise 8).

\begin{enumerate}
\def\labelenumi{\alph{enumi})}
\setcounter{enumi}{1}
\tightlist
\item
  Let \(J(\Gamma)\) be the set of subsets \(A\) of \(X\) which have the following property: for each \(x \notin A\) there is a mapping \(f \in \Gamma\) such that \(f(x) \neq x\) which leaves fixed every element of \(A\) . Show that every Hausdorff topology on \(X\) , in which every mapping \(u \in \Gamma\) is a homeomorphism of \(X\) onto itself, is finer than \(\mathcal{C}_{J}(\Gamma) = \mathcal{C}(\Gamma, J(\Gamma))\) .
\end{enumerate}

Show that the topology \(\mathfrak{G}_{\mathrm{J}}(\Gamma)\) is Hausdorff if and only if the group \(\Gamma\) satisfies the following condition: given any two distinct elements \(x, y\) of \(\mathbf{X}\) there is a finite number of elements \(u_{k} \in \Gamma\) such that, if \(\mathbf{F}(u_{k})\) is the set of elements of \(\mathbf{X}\) which are invariant under \(u_{k}\) , the \(\mathbf{F}(u_{k})\) cover \(\mathbf{X}\) and none of the \(\mathbf{F}(u_{k})\) contains both \(x\) and \(y\) .

\begin{enumerate}
\def\labelenumi{\alph{enumi})}
\setcounter{enumi}{2}
\item
  Let \(\mathcal{C}_{0}\) denote the topology of the rational line \(\mathbf{Q}\) (*resp. the real line \(\mathbf{R}_{*}\) ) and let \(\Gamma\) be the group of all homeomorphisms of \(\mathbf{Q}\) (*resp. \(\mathbf{R}_{*}\) ) onto itself. Show that the topology \(\mathcal{C}_{\mathbf{J}}(\Gamma)\) is the same as \(\mathcal{C}_{0}\) .
\item
  Let X be the subspace of the rational line consisting of 0 and the points \(\iota/n\) , where n is an integer \(\geqslant\iota\) , and let \(\Gamma\) be the group of homeomorphisms of X. Show that the topology \(\mathcal{G}_{\mathrm{J}}(\Gamma)\) is not Hausdorff, and that the group of homeomorphisms of X with respect to this topology is \(\Gamma\) .
\item
  Let X be a topological space such that for each pair of distinct points x, y of X there exists a homeomorphism u of X onto itself such that \(u(x) = x\) and \(u(y) \neq y\) (*for example, the spaces \(R^{n}_{*}\) ). Let G be the group of homeomorphisms of X onto itself, and for each \(x \in X\) let \(S_{x}\) be the subgroup of G which leaves x fixed. Show that \(S_{x}\) is equal to its normalizer in G and that the intersection of all the subgroups
\end{enumerate}

\(S_{x}\) consists of the identity element alone; in particular, the centre of G consists of the identity element alone.

\begin{enumerate}
\def\labelenumi{\arabic{enumi})}
\setcounter{enumi}{13}
\tightlist
\item
  Show that the axiom \((\mathrm{O}_{\mathrm{III}})\) is equivalent to each of the following axioms:
\end{enumerate}

\((\mathrm{O}_{\mathrm{III}}^{\prime\prime})\) If F is any closed subset of X then the intersection of all the closed neighbourhoods of F is equal to F.

\((\mathrm{O}_{\mathrm{III}}^{\prime\prime})\) If F is any filter on X which converges to a point a, then the filter F on X which has the closures of the sets of F as a base also converges to a.

\begin{enumerate}
\def\labelenumi{\arabic{enumi})}
\setcounter{enumi}{14}
\tightlist
\item
  \begin{enumerate}
  \def\labelenumii{\alph{enumii})}
  \tightlist
  \item
    Show that every Kolmogoroff space which satisfies axiom (O \(_{III}\) ) is Hausdorff (and hence regular).
  \end{enumerate}
\end{enumerate}

\begin{enumerate}
\def\labelenumi{\alph{enumi})}
\setcounter{enumi}{1}
\tightlist
\item
  Construct a non-Hausdorff topology on a set of three elements which satisfies axiom \((\mathrm{O}_{\mathrm{III}})\) .
\end{enumerate}

\begin{enumerate}
\def\labelenumi{\arabic{enumi})}
\setcounter{enumi}{15}
\item
  Let \((\mathbf{X}_t)\) be an infinite family of topological spaces, and give the product set \(\mathbf{X} = \prod_{i} \mathbf{X}_i\) one of the topologies defined in Exercise 9 of §4. Show that if the \(\mathbf{X}_t\) are Kolmogoroff (resp. accessible, Hausdorff, regular) spaces then so is \(\mathbf{X}\) .
\item
  Show that the topologies \(\mathcal{C}_0(\mathbf{X})\) , \(\mathcal{C}_+(\mathbf{X})\) and \(\mathcal{C}_{-}(\mathbf{X})\) (§ 2, Exercise 5) on a linearly ordered set \(\mathbf{X}\) are regular.
\item
  Let \(X_{1}\) , \(X_{2}\) be two sets, let \(F_{i}\) be a filter on \(X_{i}\) (i = 1, 2), and let f be a mapping of \(X_{1} \times X_{2}\) into a regular space Y such that, for each
\end{enumerate}

\[
x _ {1} \in \mathrm{X} _ {1} \text { there   exists } \lim _ {x _ {2}, \mathfrak {F} _ {2}} f (x _ {1}, x _ {2}) = g (x _ {1}).
\]

Show that a point \(a \in Y\) is a cluster point of \(g\) with respect to the filter \(\mathfrak{F}_1\) if and only if, given any neighbourhood \(V\) of \(a\) in \(Y\) and any set \(\mathbf{A}_1 \in \mathfrak{F}_1\) , there exists \(x_1 \in \mathbf{A}_1\) and \(\mathbf{A}_2 \in \mathfrak{F}_2\) such that \(f(\{x_1\} \times \mathbf{A}_2) \subset V\) . ¶ 19) Let \(X\) be a Hausdorff space which is not regular (cf.~Exercise 20), and let \(a\) be a point of \(X\) such that there is a neighbourhood \(U\) of \(a\) which does not contain any closed neighbourhood of \(a\) . Let \(\mathfrak{F}\) be the neighbourhood filter of \(a\) , and let \(Y = \mathfrak{F} \cup \{\omega\}\) be the topological space associated with the section filter of \(\mathfrak{F}\) . Let \(Z\) be the set \(Y \times X\) with the following topology: for any point \((y, z) \neq (\omega, a)\) , the neighbourhoods of \((y, z)\) are the same as for the product topology, and the neighbourhoods of \((\omega, a)\) are the sets containing sets of the form

\[
(\{\omega \} \times \mathrm{V}) \cup ((\mathrm{Y} - \{\omega \}) \times \mathrm{X}),
\]

where \(\mathbf{V}\) is a neighbourhood of \(a\) in \(\mathbf{X}\) . Let \(\mathbf{A}\) be the subset of \(\mathbf{Z}\) which is the union of the sets \(\{\mathbf{V}\} \times \mathbf{V}\) where \(\mathbf{V} \in \mathfrak{F}\) . Let \(f\) be the mapping \((V, x) \rightarrow x\) of A into X. Show that f has a limit relative to A at each point of \(\overline{A}\) , but that there is no continuous mapping of \(\overline{A}\) into X which extends f.

¶ 20) a) An open subset U of a topological space X is said to be regular if \(\mathbf{U} = \alpha(u)\) (§ 1, Exercise 3); equivalently, U is regular if U is the interior of a closed set. X is said to be semi-regular (and its topology C is said to be semi-regular) if the regular open subsets of X form a base of C. Show that if C satisfies axiom \((\mathrm{O}_{\mathrm{III}})\) then C is semi-regular.

\begin{enumerate}
\def\labelenumi{\alph{enumi})}
\setcounter{enumi}{1}
\item
  Prove that the intersection of two regular open sets is a regular open set. Deduce that the regular open sets with respect to \(\mathcal{C}\) form a base of a topology \(\mathcal{C}^*\) on X, which is coarser than \(\mathcal{C}\) and semi-regular; \(\mathcal{C}^*\) is said to be the semi-regular topology associated with \(\mathcal{C}\) . Show that for each subset U of X which is open in \(\mathcal{C}\) , the closure of U with respect to \(\mathcal{C}^*\) is the same as its closure in \(\mathcal{C}\) , and that if X is accessible for \(\mathcal{C}\) the isolated points of X with respect to \(\mathcal{C}^*\) are the same as the isolated points of X with respect to \(\mathcal{C}\) . Show that \(\mathcal{C}^*\) is Hausdorff if and only if \(\mathcal{C}\) is {[}use Exercise 3 d) of §1{]}, and finally that if Y is a regular space then the continuous mappings of X into Y are the same in both topologies \(\mathcal{C}\) and \(\mathcal{C}^*\) .
\item
  Let X be a semi-regular space and let \(C_{0}\) be its topology. Show that a topology C on X is such that \(C^{*} = C_{0}\) if and only if there is a set M of dense subsets of X (in the topology \(C_{0}\) ) such that every finite intersection of sets of M belongs to M, and such that the topology C is generated by the union of M and the set of subsets of X which are open in the topology \(C_{0}\) . {[}Consider the dense open subsets (in the topology C) and notice that every open set in C is the intersection of a dense open set (in the topology C) and an open set (in the topology \(C_{0}\) ).{]} Hence construct examples of non-regular Hausdorff topologies which are finer than a regular topology (*and even finer than a topology of a compact metrizable space*).
\end{enumerate}

\begin{enumerate}
\def\labelenumi{\arabic{enumi})}
\setcounter{enumi}{20}
\tightlist
\item
  Let X be a Hausdorff space with no isolated points, and let \(C_{0}\) be the topology of X. Let M denote the set of complements of countable subsets A of X such that \(\overline{A}\) has only a finite number of non-isolated points. Show that if C is the topology on X generated by the union of M and the set of subsets which are open in \(C_{0}\) , then \(C^{*} = C_{0}\) (Exercise 20) and that every sequence which converges with respect to C has only a finite number of distinct terms (which implies that, with respect to C, no point of X has a countable fundamental system of neighbourhoods, although X itself may be countable and therefore every point of X is then the intersection of a countable family of neighbourhoods of this point).
\end{enumerate}

¶ 22) a) With the notation of Exercise 20 c), show that the set \(\mathrm{E}(\mathfrak{C}_{0})\) of topologies C such that \(C^{*}=C_{0}\) is inductive with respect to the relation ``C is coarser than \(C'\) ''. If \(C_{0}\) is a semi-regular topology on X, a maximal element of \(\mathrm{E}(\mathfrak{C}_{0})\) is called a submaximal topology, and the set X with such a topology is called a submaximal space. Thus every topology of \(\mathrm{E}(\mathfrak{C}_{0})\) is coarser than some submaximal topology.

\begin{enumerate}
\def\labelenumi{\alph{enumi})}
\setcounter{enumi}{1}
\item
  A topology \(\mathfrak{C}\) on \(\mathbf{X}\) is submaximal if and only if every subset of \(\mathbf{X}\) which is dense in the topology \(\mathfrak{C}\) is open in \(\mathfrak{C}\) . A non-empty submaximal space is therefore not solvable (§ 2, Exercise 2).
\item
  Show that every subspace of a submaximal space is locally closed and submaximal (consider first the open subspaces and the closed subspaces).
\item
  Show that the space associated with a filter which is finer than the filter of complements of finite subsets of an infinite set is a submaximal space.
\item
  If M is a subset of a submaximal space X, show that the subspace \(\overline{M} - \mathring{M} = \operatorname{Fr}(M)\) is discrete.
\item
  Conversely, let X be a topological space which has a dense open subset U such that U is submaximal and the topology of X is U-maximal (§ 3, Exercise 11). Then X is submaximal.
\end{enumerate}

\begin{enumerate}
\def\labelenumi{\arabic{enumi})}
\setcounter{enumi}{22}
\item
  Let X be a topological space and let A be a dense subset of X such that the topology of X is A-maximal (§ 3, Exercise 11). Let f be a continuous mapping of A into a topological space Y, such that at every point of X --- A the set of cluster points of f relative to A is not empty (*this will be the case, for example, if X is not empty and Y is quasi-compact*). Show that f can be extended by continuity to the whole of X.
\item
  In the interval {]}0,1{[} of the rational line, let A be the set of rational numbers of the form \(k/2^n\) , and B the set of rational numbers of the form \(k/3^n\) ( \(k\) an integer). Let X be the space obtained by giving {]}0,1{[} the topology generated by the open intervals {]}α,β{[} contained in X, and the sets A and B. Show that X is Hausdorff and semi-regular (Exercise 20) but not regular.
\end{enumerate}

¶ 25 a) On a set X, the set of regular topologies and the set of regular topologies with no isolated points are inductive sets with respect to the relation ``C is coarser than \(C'\)''.

\begin{enumerate}
\def\labelenumi{\alph{enumi})}
\setcounter{enumi}{1}
\item
  A topology \(\mathcal{G}\) on a set X is said to be ultraregular if it is maximal in the set of topologies on X which are regular and have no isolated points; hence given any regular topology on X with no isolated points, there is a finer ultraregular topology. A space is said to be ultraregular if its topology is ultraregular. Show that a regular topology \(\mathcal{G}\) with no isolated points is ultraregular if and only if, whenever A and B are complementary subsets of X which have no isolated points, then A and B are open sets. In particular, an ultraregular space is not solvable (§ 2, Exercise 11). (To prove that the condition is sufficient, observe that if \(\mathfrak{C}'\) is a regular topology with no isolated points and is finer than \(\mathfrak{G}\) , and if U is a non-empty open set with respect to \(\mathfrak{C}'\) and \(\overline{\mathrm{U}}\) is its closure in the topology \(\mathfrak{C}'\) , then \(\overline{\mathrm{U}}\) and \(\mathrm{X} - \overline{\mathrm{U}}\) have no isolated points with respect to the topology \(\mathfrak{C}'\) ).
\item
  In an ultraregular space, the closure of every set which has no isolated points is open, and the interior of a dense set is dense; the intersection of two dense sets is therefore dense. Deduce that if \(C_{0}\) is an ultraregular topology, then amongst all the topologies C such that \(C^{*} = C_{0}\) (Exercise 20) there is one (necessarily submaximal) which is finer than all the others.
\end{enumerate}

* 26) Let \((\theta_{n})\) be a sequence of irrational numbers which is dense in the interval \(\mathbf{I} = [\mathbf{o},\mathbf{i}]\) of \(\mathbf{R}\) . Let \(\mathbf{I}_n\) be the quotient space of \(\mathbf{I}\) obtained by identifying all the points \(0, \theta_1, \ldots, \theta_n\) , and let \(\varphi_n\) be the canonical mapping \(\mathbf{I} \to \mathbf{I}_n\) . Let \(\mathbf{X}\) be the set of rational numbers contained in \(\mathbf{I}\) ; then the restriction of \(\varphi_n\) to \(\mathbf{X}\) is a bijection of \(\mathbf{X}\) onto \(\varphi_n(\mathbf{X})\) . Let \(\mathcal{G}_n\) be the inverse image under this bijection of the topology induced on \(\varphi_n(\mathbf{X})\) by the topology of \(\mathbf{I}_n\) . Show that the greatest lower bound of the decreasing sequence \((\mathcal{G}_n)\) of Hausdorff topologies on \(\mathbf{X}\) is not a Hausdorff topology.*

¶ 27) Let X be a topological space. Show that there exists a Kolmogoroff (resp. accessible, Hausdorff, regular) space Y and a continuous mapping \(\varphi: X \to Y\) which has the following universal property: for every continuous mapping \(f: X \to Z\) of X into a Kolmogoroff (resp. accessible, Hausdorff, regular) space Z, there is a unique continuous mapping \(g: Y \to Z\) such that \(f = g \circ \varphi\) (Set Theory, Chapter IV, § 3, no. 1). Prove that:

\begin{enumerate}
\def\labelenumi{\alph{enumi})}
\item
  For Kolmogoroff spaces, we may take Y to be the quotient space of X by the equivalence relation \(\{\bar{x}\}=\{\bar{y}\}\) .
\item
  For accessible spaces, we may take Y to be the quotient space of X by the following equivalence relation R: the graph of R is the intersection of all graphs G of equivalences on X such that \(\mathrm{G}(x)\) is closed in x for each \(x \in X\) ; the graph of R then contains the graph of the relation \(\{\bar{x}\} \cap \{\bar{y}\} \neq \emptyset\) between x and y.
\item
  For Hausdorff spaces and regular spaces, show that conditions \((\mathrm{CU}_{\mathbf{I}})\) to \((\mathrm{CU}_{\mathbf{III}})\) of Set Theory, Chapter IV, § 3, no. 2, are satisfied, by using in particular Exercise 6 a). Give examples in which X is not regular and the mapping \(\varphi: X \to Y\) of X into the corresponding universal regular space is bijective {[}cf.~Exercise 20 b){]}.
\end{enumerate}

\begin{enumerate}
\def\labelenumi{\alph{enumi})}
\setcounter{enumi}{3}
\tightlist
\item
  If X is a topological space which satisfies axiom \((\mathrm{O}_{\mathrm{III}})\) , show that the universal Kolmogoroff space corresponding to X is canonically homeomorphic to the universal regular space corresponding to X.
\end{enumerate}

\begin{enumerate}
\def\labelenumi{\arabic{enumi})}
\setcounter{enumi}{27}
\tightlist
\item
  Let X be a Hausdorff space which is not regular, let F be a closed subset of X and let \(a \notin F\) be a point of X such that every neighbourhood of a meets every neighbourhood of F. Let R be the equivalence relation on X obtained by identifying all the points of F. Show that R is closed and that the graph of R is closed in \(X \times X\) , but that R is not Hausdorff.
\end{enumerate}

¶ 29) a) Let R be a Hausdorff equivalence relation on a non-empty topological space X, and let \(\mathfrak{F}(X)\) be the space of non-empty closed subsets of X, with the topology induced by \(\mathfrak{G}_{\Theta}\) (Exercise 12). Show that every \(A \in \mathfrak{F}(X)\) which lies in the closure of X/R {[}considered as a subset of \(\mathfrak{F}(X)\) {]} is contained in some equivalence class of R (use reductio ad absurdum).

\begin{enumerate}
\def\labelenumi{\alph{enumi})}
\setcounter{enumi}{1}
\tightlist
\item
  Suppose in addition that X is regular and R is open. Show that X/R is then a closed subset of \(\mathfrak{F}(X)\) with respect to the topology \(\mathfrak{G}_{\Theta}\) (use Proposition 11 of § 7, no. 6).
\end{enumerate}

